\documentclass[conference]{IEEEtran}

\usepackage[utf8]{inputenc}
\usepackage[T1]{fontenc}

\usepackage{hyperref}
\usepackage{listings}
\usepackage{listingsutf8}
\usepackage{setspace}
\usepackage{graphicx}
\usepackage{amsmath}

\lstset{
  inputencoding=utf8,
  extendedchars=true,
  basicstyle=\ttfamily\footnotesize,
  breaklines=true,
  breakatwhitespace=false,
  columns=fullflexible,
  keepspaces=true,
  showstringspaces=false,
  literate=
    {’}{{'}}1   % 스마트 아포스트로피 → '
    {‘}{{`}}1   % 여는 따옴표 → `
    {“}{{``}}1  % 여는 큰따옴표 → ``
    {”}{{''}}1  % 닫는 큰따옴표 → ''
    {–}{{-}}1   % en dash → -
    {—}{{-}}1   % em dash → -
}

\begin{document}

\title{Speaker-Adaptive Voice Assistant for Multi-User Homes}

\author{
\IEEEauthorblockN{
Kun Lee,
Jinseo Hong,
Patrick Segedi
}
\IEEEauthorblockA{
Department of Information Systems, Hanyang University \\
Department of Computer Science, Chalmers University of Technology \\
Emails: ceh1502@hanyang.ac.kr, h0dduck@hanyang.ac.kr, segedi@chalmers.se
}
}

\maketitle

\begin{abstract}
This document contains the full content of the README.md project blog for the G08 AI and Applications Project, reformatted into the IEEE conference template without modifying the original text.
\end{abstract}

\section*{README Content (Unmodified)}

\begin{lstlisting}
G08-ai-and-applications-project
Title: Speaker-Adaptive Voice Assistant for Multi-User Homes
blog: https://patricksegedi.github.io/G08-ai-and-applications-project

Members: 

Kun Lee, Dept. of Information Systems Hanyang University, ceh1502@hanyang.ac.kr

Jinseo Hong, Dept. of Information Systems Hanyang University, h0dduck@hanyang.ac.kr

Patrick Segedi, Dept. of Computer Science Chalmers Unviersity of segedi@chalmers.se

### I. Introduction
- Motivation: Why are you doing this?
Our team embarked on this project with two primary motivations. 
First, we wanted to gain hands-on experience applying cutting-edge 
AI technologies to solve real-world problems. Voice recognition 
and speaker verification represent some of the most exciting 
frontiers in artificial intelligence, and we saw this as an 
invaluable opportunity to deepen our understanding of these 
technologies beyond theoretical knowledge.

Second, this project brought together students from three different 
universities across three countries - Hanyang University (Korea), 
Chalmers University of Technology (Sweden), and Zurich University 
of Applied Sciences (Switzerland). This international collaboration 
allowed us to experience diverse perspectives, work across time zones, 
and develop the communication skills essential for global tech careers.

- What do you want to see at the end?
Our ultimate goal is to see our voice assistant actually working - 
responding to real voices, recognizing different family members, 
and controlling smart home devices accordingly. We want to demonstrate 
that a privacy-preserving, speaker-adaptive voice assistant is not 
just a theoretical concept, but a functional system that can enhance 
daily life in multi-user households.

### II. Datasets

Our project uses three categories of datasets:
1. pretrained public datasets used by underlying AI models,
2. user-generated enrollment voice data, and
3. real-time audio streams captured by in-home devices.
These datasets support speaker verification, speech recognition, and contextual inference. 

a) Pretrained Speaker Verification Dataset (VoxCeleb)
The ECAPA-TDNN model from SpeechBrain, used for speaker identification, is pretrained on the VoxCeleb1/2 dataset, a large-scale speaker recognition corpus containing thousands of speakers recorded in unconstrained, real-world environments.

- Purpose in our system: extracting robust speaker embeddings for multi-user identity recognition.
We do not retrain the model but rely on the pretrained embeddings to authenticate speakers locally on the device.

b) Pretrained ASR Dataset (Whisper / Faster-Whisper)
Our wake-word detection and speech-to-text pipeline uses Faster-Whisper, an optimized implementation of OpenAI’s Whisper model. Whisper is trained on:

- 680,000 hours of multilingual speech

- Noisy, real-world audio data

- Diverse accents, environments, and speaking styles
This makes it suitable for household environments where commands may vary in tone or clarity.

c) User Enrollment Voice Dataset (Locally Stored)
To enable personalization and role-based access control, each household member records a set of enrollment samples during profile setup.

- Format: WAV, 16 kHz, mono

- Typical length: 3–5 seconds per sample

- Number of samples: 3–5 recordings per user

- Storage: encrypted local SQLite database

- Purpose:

  - create user-specific embeddings

  - authenticate speakers in real-time

  - maintain privacy by preventing cloud transmission
    
All data remains fully on-device to ensure privacy.

d) Real-Time Environmental Audio Stream
For wake-word detection, command processing, and location awareness, the system processes continuous audio input from multiple room speakers/microphones.
- Purpose:

  - detect speaker location

  - identify the active user

  - resolve multi-user command conflicts
  - 
Since this stream is processed on-device and never uploaded, it is considered an operational dataset rather than a stored one.

### III. Methodology 

Faster-Whisper:
OpenAI’s speech-to-text model in an optimized CTranslate2 implementation. Used to convert user speech into text for wake-word detection and command processing.

SpeechBrain (ECAPA-TDNN):
Used for speaker verification. The ECAPA-TDNN model generates speaker embeddings to compare prerecorded samples with user speech.

Cosine Similarity Scoring:
Similarity between voice embeddings is computed using cosine similarity to identify the best-matching speaker.

Wake-Word Detection (Keyword-based):
A lightweight text-based wake-word detector checks whether the transcribed text contains the activation word (“Hello”).

XTTS-v2 (Coqui TTS):
A neural TTS model used to generate spoken responses with selected speaker voices.

Audio Capture (SoundDevice + WAVIO):
Used to record raw microphone audio (16 kHz mono WAV format) as input for STT and verification.

Lovelace UI:
Home Assistant’s dashboard system, used to visualize interactions with IoT devices.

### IV. Evaluation & Analysis
none

### V. Related Work (e.g., existing studies)
none

### VI. Conclusion: Discussion
Coordinating schedules was still difficult due to different class times and personal commitments. 


We initially believed that adopting Agile with frequent meetings would 
improve efficiency, but we learned that meeting more often does not 
necessarily lead to better productivity—what mattered more was the 
quality of communication and clear task ownership between meetings.

## Weekly progress 
### Week 1 - Project Start [2025-10-12]

During the first week, we set up GitHub, platforms, planned team meetings, and assigned responsibilities.

- Created a GitHub repo.
- Set up VS Code.
- Planned group meetings and roles.

Next steps:
- Research project ideas.

Summary: We established the project structure and clarified team responsibilities.

### Week 2 - Research project ideas [2025-10-19]

During the second week, each group member explored and compared potential ideas for the project.
- Each group member explored ideas for the school project.

Next steps:
- Review project ideas together.
- Choose a project.

Summary: We gathered individual ideas and prepared to decide on a final project direction.

### Week 3 - Choose a project [2025-10-26]

During the third week, we discussed our ideas together and selected a final project direction.
- Discussed ideas together and picked a final direction.
- Chosen project: Context-aware Personalized Voice Assistant.

Next steps:
- Research pre-trained voice identification models.
- Explore implementation frameworks and tools.

Summary: We agreed on the final project concept and began investigating suitable models and technologies.

### Week 4 -  Model exploration [2025-11-02]

During the fourth week, we reviewed available technologies and explored potential approaches for our system.
- Reviewed findings, potential technologies, and tools together.
- Decided to use the pre-trained ECAPA-TDNN model from the SpeechBrain toolkit.
- Planned implementation in Python.

Next steps:
- Individually test the model

Summary: We selected the model and defined the technical direction for implementation.

### Week 5 -  Architecture & prototype [2025-11-09]

During the fifth week, we focused on the system architecture and began building the first version of our project.
- Researched architecture and design approaches.
- Clarified roles and implementation responsibilities and divided tasks.
- Built the first prototype.
- Implemented Agile workflow using a Scrum/Kanban board.

Next steps:
- Begin database implementation.
- Develop a prototype using the Lovelace smart-home simulation environment to demonstrate the voice assistant in the final presentation.
- Continue testing and refining the prototype.

Summary: We established the system architecture, created the first prototype, and organized the development workflow using Agile practices.

### Week 6 - Implementation [2025-11-16]

During the six week, we continued working with implementation of the different parts of the program.
- Finished the LoveLace smart-home simulation.
- 0Continuted with database implementaton.
- Improved the prototype and finished the wake-word implementation and also demostrated it to professor.

Next steps:
- Documentation managment
- Create a more detailed architecture visualisation.
- Completed the database implementation.

Summary: Completed the wake-word implementation and LoveLace home simulation. We also demostrated our progress to professor and received feedback on our current work.

### Week 7 - Documentation [2025-11-23]

During the seven week, we mostly focused on documentation for the project, while also continued to work on the database and started working on the back-end implemetation.
- The main focus was documentation for the project.
- Continuted with database implementaton.
- Back-end development.
- Updated the document from docx-format to LateX.
- Architecture visualization
- Presentation discussion.

Next steps:
- Finishing the database
- Continue working on documentation. Adding use cases, model description, etc.
- Model improvement
- Web application development
- Organize the software dictionary.
- Presentation planning.

Summary: This week we focused on improving the project documentation after receiving feedback from our latest meeting with Professor. We learned how to use LateX and updated the documentation from .docs to LateX. We also made a program architecture visualization map. The database implementation is currently being worked on and back-end implementation has started.

### Week 8 - updating implementation [2025-11-30]

During Week 8, we focused on connecting all core components through a fully structured backend.
We built a FastAPI server, implemented API endpoints, and enabled automatic documentation using Swagger UI.
- we constructed a MySQL database and configured the backend to interact with it—handling user information
- device data
- all required project entities
- ensuring that the backend, database, and voice-analysis modules could communicate seamlessly in a unified system

Next steps:
- Finishing the implementation
- Testing
- Creating the demo video

### Week 9 - Testing & making Demo Video [2025-12-07]

During week 9 after confirming stable performance across all modules, we recorded the official demo video.
- Filming our actual demonstration
- Merging into visual video
- Testing the demo with no errors.

## Task 2 - Video URL

https://www.youtube.com/watch?v=uW7JgBTDY0g

\end{lstlisting}

\end{document}
